% this document uses paper size A4, default fontsize 12 and properties from the KoMa script scartl class
\documentclass[a4paper,12pt]{scrartcl}

% these packages are often needed but not strictly necessary
%\usepackage{numprint}
%\usepackage{pgfplots}
%\pgfplotsset{compat=1.16}
%\usepackage{amsmath}
%\usepackage{pgffor}
%\usepackage[roman]{parnotes}
%\usepackage{xifthen}

% necessary imports
\usepackage[ngerman]{babel}
\usepackage[]{csquotes}
\usepackage{titling}
\usepackage{fancyhdr}
\usepackage{graphicx}
\usepackage{float}
\usepackage[margin=2.5cm,includeheadfoot, footskip=80pt]{geometry}
\usepackage[breaklinks]{hyperref}
\usepackage{blindtext}
\usepackage[nospace]{varioref}
\usepackage[backend=biber, style=apa,sorting=nyt]{biblatex}
\usepackage{booktabs}
\usepackage[justification=centering]{caption}
\usepackage[titles]{tocloft}
\usepackage{tabularx}
\usepackage{setspace}
\usepackage{sectsty}
\usepackage{units}
\usepackage{xcolor}
\usepackage{parskip}
\usepackage{hyperref}
\usepackage{listings}
\usepackage{color}

%\definecolor{darkgray}{gray}{0.2} % Dunkelgrauer Hintergrund
%\definecolor{dkgreen}{rgb}{0,0.6,0}
%\definecolor{gray}{rgb}{0.5,0.5,0.5}
%\definecolor{mauve}{rgb}{0.58,0,0.82}
%
%\lstset{ 
%  language=python,
%  aboveskip=3mm,
%  belowskip=3mm,
%  showstringspaces=false,
%  columns=flexible,
%  basicstyle={\small\ttfamily\color{white}}, % Textfarbe weiß
%  backgroundcolor=\color{darkgray},          % Dunkler Hintergrund
%  numbers=left,
%  numberstyle=\tiny\color{gray},
%  keywordstyle=\color{blue},
%  commentstyle=\color{dkgreen},
%  stringstyle=\color{mauve},
%  breaklines=true,
%  breakatwhitespace=true,
%  tabsize=3,
%  frame=none,  % Keine Linien oben und unten
%  lineskip=-0.5pt
%}

%Light Mode
\definecolor{dkgreen}{rgb}{0,0.6,0}
\definecolor{gray}{rgb}{0.5,0.5,0.5}
\definecolor{mauve}{rgb}{0.58,0,0.82}

\lstset{frame=tb,
  language=python,
  aboveskip=3mm,
  belowskip=3mm,
  showstringspaces=false,
  columns=flexible,
  basicstyle={\small\ttfamily},
  numbers=left,
  numberstyle=\tiny\color{gray},
  keywordstyle=\color{blue},
  commentstyle=\color{dkgreen},
  stringstyle=\color{mauve},
  breaklines=true,
  breakatwhitespace=true,
  tabsize=3
}

% set metadata for the PDF
\hypersetup{
    pdfauthor={Your Name},      % Add your name here
    pdftitle={The Title},       % Add the document title here
    pdfsubject={The Subject},   % Add the subject of the document here
    pdfkeywords={Some Keywords} % Add relevant keywords here
}

% location of bib and graphics
\addbibresource{myBib.bib}
\graphicspath{{images/}}

% definitions

\definecolor{ba-blau}{HTML}{093a80}
\newcommand{\documentTitle}[1]{{\fontsize{20pt}{24pt}\selectfont\textbf{\color{ba-blau}#1}}}
\newcommand{\pageTitle}[1]{{\fontsize{14pt}{17pt}\selectfont\textbf{\color{ba-blau}#1}}}
% define command for for using full reference in footnote
\newcommand{\fullfootcite}[1]{
  \footnote{\fullcite{#1}}
}

% styling

% set default line spacing
\setstretch{1.5}

% set section headings fontsizes and line spacing
\sectionfont{\fontsize{14}{16.8}\selectfont\fontfamily{ptm}\selectfont}
\subsectionfont{\fontsize{12}{14.4}\selectfont\fontfamily{ptm}\selectfont}
% color section headings
\addtokomafont{disposition}{\color{ba-blau}}
\addtokomafont{sectionentry}{\color{ba-blau}}

% styling of the title page header and footer
\fancypagestyle{Deckblatt}{
    \setlength{\headheight}{10pt}
    \pagestyle{fancy}
    \fancyhead{}
    \setlength{\headheight}{2cm}
    \setlength{\headsep}{80pt}
    \chead{\includegraphics[]{deckbild.jpeg} \\}
    %pagefooter
    \fancyfoot{}
    \fancyfoot[R]{\thepage}
}
% styling of the default page header and footer
\fancypagestyle{defaultPageStyle}{
    \setlength{\headheight}{10pt}
    \pagestyle{fancy}
    \fancyhead{}
    \setlength{\headheight}{2cm}
    \setlength{\headsep}{40pt}
    \chead{\includegraphics[]{deckbild.jpeg} \\}
    %pagefooter
    \fancyfoot{}
    \fancyfoot[R]{\thepage}
}

% start of document

\begin{document}

% use different pagestyling
\thispagestyle{Deckblatt}
% ignore this page when numbering
\pagenumbering{gobble}

% title page
\begin{center}
    \documentTitle {Titel der Bachelor Thesis} \\[3ex]
    \includegraphics[width=3cm]{Platzhalter_logo.png} \\[3ex]
    \large{Exposé erstellt im Rahmen des Bachelor Thesis Kolloquiums:
    } \\
    \large{DD.MM.YYYY}
\end{center}

\renewcommand{\arraystretch}{2.5}
\begin{table}[h]
    \setlength{\tabcolsep}{32pt}
    \begin{tabularx}{\textwidth}{l l}
        Studiengruppe:                      & XX\_XS2X\_XX     \\
        Name Studierender:                  & Max Mustermann   \\
        Anzahl der Wörter:                  & 0                \\ [-15pt]
        (inkl. wörtliche Zitate / Fußnoten) & ~                \\
        Anzahl der Wörter:                  & 0                \\ [-15pt]
        (exkl. wörtliche Zitate / Fußnoten) & ~                \\
        Akademischer Gutachter:             & Prof. Mustermann \\
        Betrieblicher Gutachter:            & Herr Firma       \\
        Abgabedatum:                        & DD.MM.YYYY       \\
    \end{tabularx}
\end{table}

\newpage

% use default page styling from now on
\pagestyle{defaultPageStyle}
\pageTitle{Sperrvermerk}
Die vorliegende Bachelor Thesis beinhaltet interne vertrauliche Informationen der XYZ GmbH/AG. Die Weitergabe des Inhaltes dieser Arbeit und eventuell beiliegender Abbildungen, Tabellen und Daten im gesamten oder in Teilen ist grundsätzlich untersagt. Es dürfen keinerlei Kopien oder Abschriften, auch nicht in digitaler Form, gefertigt werden. Ausnahmen bedürfen der schriftlichen Genehmigung durch die XYZ GmbH/AG. 
\newpage
\pageTitle{Gleichbehandlung der Geschlechter}\\
Aus Gründen der besseren Lesbarkeit wird bei Personenbezeichnungen und personenbezogenen Hauptwörtern in diesem Dokument die maskuline Form verwendet. Entsprechende Begriffe gelten im Sinne der Gleichbehandlung grundsätzlich für alle Geschlechter. Die verkürzte Sprachform hat nur redaktionelle Gründe und beinhaltet keine Wertung.
\newpage
% Wird ein solcher Hinweis nicht vorangestellt, muss laut den Vorgaben geschlechterinklusive Sprache in der Arbeit genutzt werden.

% start page numbering in roman numerals
\pagenumbering{Roman}
\setcounter{page}{1}

\tableofcontents
\newpage
\listoffigures
\newpage
\listoftables
\newpage

% store last roman page number
\newcounter{preamblecounterstate}
\setcounter{preamblecounterstate}{\value{page}}

% use arabic numbers for actual content pages
\pagenumbering{arabic}

% include the pages stored in the section folder
\section{Demo-Seite}
%Falls keine Seitenanzahl vorhanden ist  dann:\footcite[Vgl.][]{DemoQuelle}
Demo-Quelle \footcite[Vgl.][\printfield{pages}]{DemoQuelle}

\begin{center}
    \begin{table}[h]
    \centering
    \begin{tabular}{|c|p{6cm}|}
        \hline
        \textbf{Datum} & \textbf{Aktivitäten} \\
        \hline
        Kebab & 7€ \\
        \hline
        Adana & \begin{itemize}
            \item \textbf{Groß}: 8€
            \item \textbf{Klein}: 6€
        \end{itemize} \\
        \hline
        Köfte & \begin{itemize}
            \item 5 Stück: 8€
            \item 2 Stück: 6€
        \end{itemize}\\
        \hline
        Mercimek Suppe & 3€ \\
        \hline
        Dönerteller & 15€ \\
        \hline
    \end{tabular}
    \caption{Beispiel Tabelle.}
    \label{tab:example}
    \end{table}
    \end{center}

Die Tabelle zeigt den Preis eines Dönertellers, dieser lässt sich wie folgt berechnen:
\begin{equation}
    15 = \sum_{n=1}^{10} \frac{n}{20} + \sum_{k=1}^{5} \frac{2k}{10} - \sum_{i=1}^{3} i
\end{equation}
\newpage
Der folgende Abschnitt könnte hilfreich für eine Ausarbeitung in der Informatik sein.
\begin{lstlisting}
    def hello_world():
        print("Hello, World!")
\end{lstlisting}
\section{Kontext der Ausarbeitung}


\subsection{Relevanz der Untersuchung}


\subsection{Stand der Forschung}


\section{Erstellung der Arbeit}
Test\fullfootcite{goos_single_2002}
\subsection{Genutzte Quellen}
\subsection{Zeitplan}
\section{Zielsetzung}


\subsection{Grundlagen}


\subsection{Forschungsfrage}



\newpage

% set page numbering to roman for bibliography
\pagenumbering{Roman}
% read counter stored earlier
\setcounter{page}{\value{preamblecounterstate}}

\printbibliography[title={Literaturverzeichnis}]

\end{document}